 %%%%%%%%%%%%%%%%%%%%%%%%%%%%%%%%%%%%%%%%%
% Beamer Presentation
% LaTeX Template
% Version 1.0 (10/11/12)
%
% This template has been downloaded from:
% http://www.LaTeXTemplates.com
%
% License:
% CC BY-NC-SA 3.0 (http://creativecommons.org/licenses/by-nc-sa/3.0/)
%
%%%%%%%%%%%%%%%%%%%%%%%%%%%%%%%%%%%%%%%%%

%----------------------------------------------------------------------------------------
%	PACKAGES AND THEMES
%----------------------------------------------------------------------------------------

\documentclass[9pt,aspectratio=169]{beamer}
\usepackage[utf8]{inputenc}
\usepackage{bm}
\usepackage{bbm}
\usepackage{graphicx}
\usepackage{cmbright}
\usefonttheme{professionalfonts}
%\usefonttheme{serif}
% The Beamer class comes with a number of default slide themes
% which change the colors and layouts of slides. Below this is a list
% of all the themes, uncomment each in turn to see what they look like.

%\usetheme{default}
%\usetheme{AnnArbor}
% \usetheme{Antibes}
%\usetheme{Bergen}
%\usetheme{Berkeley}
%\usetheme{Berlin}
%\usetheme{Boadilla}
% \usetheme{CambridgeUS}
%\usetheme{Copenhagen}
%\usetheme{Darmstadt}
% \usetheme{Dresden}
\usetheme{Frankfurt}
%\usetheme{Goettingen}
% \usetheme{Hannover}
% \usetheme{Ilmenau}
% \usetheme{JuanLesPins}
%\usetheme{Luebeck}
%\usetheme{Madrid}
%\usetheme{Malmoe}
%\usetheme{Marburg}
% \usetheme{Montpellier}
% \usetheme[width=2.5cm]{PaloAlto}
%\usetheme{Pittsburgh}
% \usetheme{Rochester}
%\usetheme{Singapore}
% \usetheme{Szeged}
% \usetheme{Warsaw}
% \setbeamertemplate{title page}[default][colsep=-4bp,rounded=false]
\setbeamertemplate{headline}{}
\defbeamertemplate{footline}{myframe number}
{
  \hspace{1pt}
    \usebeamercolor[structure]{page number in head/foot}%
  \usebeamerfont{page number in head/foot}%
  \raisebox{1.7ex}[0pt][0pt]{% <--- change here
    \insertframenumber\,/\,\inserttotalframenumber\kern1em}%
}
\setbeamertemplate{footline}[myframe number]
\AtBeginSection[]
{
  \begin{frame}
    \frametitle{Outline}
    \begin{minipage}{1.0\linewidth}
      \tableofcontents[currentsection,currentsubsection]
    \end{minipage}
  \end{frame}
}

\AtBeginSubsection[]
{
  \begin{frame}
    \frametitle{Outline}
    \begin{minipage}{1.0\linewidth}
      \tableofcontents[currentsection,currentsubsection]
    \end{minipage}
  \end{frame}
}

% As well as themes, the Beamer class has a number of color themes
% for any slide theme. Uncomment each of these in turn to see how it
% changes the colors of your current slide theme.
%\definecolor{mycolor}{rgb}{.8,.0,.3}
%\setbeamercolor*{palette primary}{use=structure,fg=white,bg=mycolor}
%\usecolortheme{albatross}
%\usecolortheme{beaver}
%\usecolortheme{beetle}
%\usecolortheme{crane}
%\usecolortheme{dolphin}
%\usecolortheme{dove}
%\usecolortheme{fly}
%\usecolortheme{lily}
%\usecolortheme{orchid}
% \usecolortheme{rose}
%\usecolortheme{seagull}
\usecolortheme{seahorse}
%\usecolortheme{whale}
%\usecolortheme{wolverine}
%\usecolortheme[named=mycolor]{structure}
%\setbeamertemplate{footline} % To remove the footer line in all slides uncomment this line
%\setbeamertemplate{footline}[page number] % To replace the footer line in all slides with a simple slide count uncomment this line

%\setbeamertemplate{navigation symbols}{} % To remove the navigation symbols from the bottom of all slides uncomment this line


\usepackage{booktabs} % Allows the use of \toprule, \midrule and \bottomrule in tables

%----------------------------------------------------------------------------------------
%	TITLE PAGE
%----------------------------------------------------------------------------------------

\title[Lecture 4: Constrained Optimization I]{Lecture 4: Constrained
  Optimization I}

\author{Junnan Zhang} % Your name
\institute[Paula and Gregory Chow Institute for Studies in Economics\\
Xiamen University] % Your institution as it will appear on the bottom of every slide, may be shorthand to save space
{
	Paula and Gregory Chow Institute for Studies in Economics \\ Xiamen University  \\ % Your institution for the title page
	\medskip
    \textit{Slides Prepared by Xiaoling Mei}
	%\textit{luciamay@smith.com} % Your email address
}
\date{Fall, 2026} % Date, can be changed to a custom date
\setbeamertemplate{itemize items}[default]
\setbeamertemplate{sidebar}[sidebar right]
% \definecolor{myco}{HTML}{a67678}
\definecolor{myco}{HTML}{8776a6}
\setbeamercolor{itemize item}{fg=myco} % all frames will have red bullets
\setbeamercolor{itemize subitem}{fg=myco} % all frames will have red bullets
\setbeamercolor{block title}{bg=myco}
\setbeamercolor{structure}{fg=myco}
% \setbeamercovered{transparent}
\newcommand{\RR}{\mathbbm R}

\begin{document}



\begin{frame}
\titlepage % Print the title page as the first slide
\end{frame}

\section{Definition}
\begin{frame}
    \frametitle{Motivation}

    We often work with the optimization problem in which the objects are not
    free to take on any value but are constrained. For example:
    \begin{itemize}
        \item A household's consumption is constrained by available income
        \item A firm's production is constrained by the cost and availability
        of its inputs
        \item The central mathematical problem here is that of
        maximizing/minimizing a function of several variables, where these
        variables are bound by some constraining equations.
    \end{itemize}
\end{frame}		
		

\begin{frame}
	\frametitle{Definitions}
	\begin{block}{\textbf{The Prototype Problem}}
	\begin{itemize}
		\item Objective function: $f(x_1,x_2,\cdots,x_n)$
		\item Constraint functions: 
		\begin{align}
		g_1(x_1,\cdots,x_n) &{\color{red}\leq} b_1 \nonumber \\
		& \vdots \nonumber \\
		g_k(x_1,\cdots,x_n) &{\color{red}{\leq}} b_k \nonumber \\ 
		h_1(x_1,\cdots,x_n) &{\color{red} =} c_1 \nonumber \\
		& \vdots \nonumber \\
		h_m(x_1,\cdots,x_n) &{\color{red} =} c_m \nonumber 
		\end{align}
	\end{itemize}
	\end{block}
\end{frame}	














\section{Equality Constraints}

\begin{frame}
	\frametitle{Equality Constraints}
\begin{itemize}
	\item Simplest constrained maximization problem: maximizing a function of two variables subject to a single equality constraint
	\item Setup: 
	\begin{align}
	\text{max}\ \ & f(x_1,x_2) \nonumber \\
	\text{s.t.}\ \ & h(x_1,x_2) = c \nonumber 
	\end{align}
	\item The highest level set of $f$ must touch (be tangent to) the
    constraint curve $C$ at the constrained max. Figure 18.1.
\end{itemize}
\end{frame}

\begin{frame}
	\frametitle{Equality Constraints}
	\begin{itemize}
		\item Equivalently, at the constrained max $\mathbf{x^*}$, the slope
        of the level set of $f$ equals the slope of the constraint curve $C$. 
		\item If $\frac{\partial f}{\partial x_2}(\mathbf{x^*})\neq0$, the
        slope of the level set of $f$ at $\mathbf{x^*}$ is
        $$-\frac{\partial f}{\partial x_1}(\mathbf{x^*})\Big/\frac{\partial f}{\partial x_2}(\mathbf{x^*})$$
		\item If $\frac{\partial h}{\partial x_2}(\mathbf{x^*})\neq0$, the
        slope of the constraint set $h(x_1,x_2) = c$ at $\mathbf{x^*}$ is
		$$-\frac{\partial h}{\partial x_1}(\mathbf{x^*})\Big/\frac{\partial h}{\partial x_2}(\mathbf{x^*})$$
	\end{itemize}
\end{frame}

	
\begin{frame}
	\frametitle{Equality Constraints}
	\begin{itemize}
        \item Equalize the slope and we get:
		$$\frac{\frac{\partial f}{\partial x_1}(\mathbf{x^*})}{\frac{\partial h}{\partial x_1}(\mathbf{x^*})} = \frac{\frac{\partial f}{\partial x_2}(\mathbf{x^*})}{\frac{\partial h}{\partial x_2}(\mathbf{x^*})} = \mu.$$
		\item Rewrite the equation and combine with the constraint equation:
		\begin{align}
		\frac{\partial f}{\partial x_1}(\mathbf{x^*})-\mu\frac{\partial h}{\partial x_1}(\mathbf{x^*}) &= 0 \nonumber \\
		\frac{\partial f}{\partial x_2}(\mathbf{x^*})-\mu\frac{\partial h}{\partial x_2}(\mathbf{x^*}) &= 0 \nonumber \\
		h(x_1,x_2) &= c \nonumber 
		\end{align}
		 A system with three unknowns.
         \item More generally, tangency means that the gradients are
         proportional:
         $$\nabla f(\mathbf{x^*})=\mu\nabla h(\mathbf{x^*}).$$

	\end{itemize}
\end{frame}


\begin{frame}
	\frametitle{Equality Constraints}
	\begin{block}{\textbf{Constraint Qualification}}
		\begin{itemize}
			\item Above method would not have worked if both $\partial
            h/\partial x_1$ and $\partial h/\partial x_2$ were 0;
			\item We will need to make the assumption that $\partial
            h/\partial x_1$ or $\partial h/\partial x_2$ (or both) is not
            zero at the maximizer
			\item This restriction is called a \textbf{constraint qualification}
		\end{itemize}
	\end{block}
\end{frame}
				
\begin{frame}
	\frametitle{Equality Constraints: Theorem}
	\begin{block}{\textbf{Theorem 18.1}}
        Suppose $\mathbf{x^*} = (x_1^*,x_2^*)$ is a solution and suppose
        further that it is not a critical point of $h$. Then there is real
        number $\mu^*$ such that $(x_1^*,x_2^*,\mu^*)$ is a critical point of
        the \emph{Lagrangian function}
	$$L(x_1,x_2,\mu) \equiv f(x_1,x_2)-\mu(h(x_1,x_2)-c).$$ In other words, at $(x_1^*,x_2^*,\mu^*)$:
	$$\frac{\partial L}{\partial x_1}=0, \frac{\partial L}{\partial x_2}=0, \frac{\partial L}{\partial \mu}=0$$
	$\mu$ is called a \textbf{Lagrange multiplier}.
	\end{block}
\end{frame}
		
						
\begin{frame}
	\frametitle{Example 18.4} \textbf{Example 18.4}: Use Theorem 18.1 to
    solve a simple utility maximization problem:
	\begin{align}
	\text{maximize} \quad f(x_1,x_2) &= x_1x_2 \nonumber \\
	\text{subject to}\quad h(x_1,x_2)&\equiv x_1+4x_2 = 16 \nonumber 
	\end{align}
    The gradient of $h$ is $(1, 4)$, so $h$ has no critical points and the
    constraint qualification is satisfied. From the Lagrangian function
$$L(x_1,x_2,\mu) = x_1x_2-\mu(x_1+4x_2-16)$$
\end{frame}			
		
		
\begin{frame}
	\frametitle{Example 18.4}
	\textbf{Example 18.4} : Set the partial derivatives equal to zero:
	\begin{align}
        \frac{\partial L}{\partial x_1} & = x_2-\mu = 0 \nonumber \\
        \frac{\partial L}{\partial x_2} & = x_1-4\mu = 0 \nonumber \\
        \frac{\partial L}{\partial \mu} & = -(x_1+4x_2-16) = 0\nonumber
    \end{align}
    We conclude the solution of this system is
    $$x_1 = 8,\, x_2 = 2,\, \mu = 2$$ So the only candidate for
    a solution is $$x_1 = 8,\, x_2 = 2$$
\end{frame}		
		
		

	


	



	


\begin{frame}
	\frametitle{Several Equality Constraints}
\begin{itemize}
	\item When there are several equality constraints, the problem becomes
	\begin{align}
	\max & \ f(x_1,x_2,\cdots,x_n) \nonumber \\
	\text{s.t.}\ \ & h_1(\mathbf{x}) = a_1; \cdots; h_m(\mathbf{x}) = a_m \nonumber 
	\end{align}
    \item Generalized constraint qualification: $x$ is called a critical
    point of $\mathbf{h} = (h_1,h_2,\cdots,h_m)$ if the rank of the matrix
    $D\mathbf h(\mathbf{x}^*)$ is less than $m$.
    \item \textbf{Nondegenerate Constraint Qualification}: The rank of
    Jacobian $D\mathbf h(\mathbf{x}^*)$ is equal to the number of the constraints,
    then $\mathbf{x^*}$ satisfies NDCQ.
\end{itemize}
\end{frame}

\begin{frame}
	\frametitle{Equality Constraints: Theorem}
	\begin{block}{\textbf{Theorem 18.2}}
		Suppose $x^*$ is a solution and suppose further it satisfies
        condition NDCQ. Then there exists $\mu^* = (\mu_1^*, \cdots,
        \mu_m^*)$ such that $(\mathbf x^*,\mu^*)$ is a critical
        point of the \emph{Lagrangian} \small
		\begin{equation*}
            L(x,\mu) \equiv f(x)-\mu_1(h_1(x)-a_1)-\mu_2(h_2(x)-a_2)
            - \cdots - \mu_m(h_m(x)-a_m).
		\end{equation*}
	 In other words, 
		\begin{align}
            \frac{\partial L}{\partial x_1}(x^*,\mu^*) &=0,
            \cdots, \frac{\partial L}{\partial
              x_n}(x^*,\mu^*)=0 \nonumber \\ 
            \frac{\partial L}{\partial \mu_1}(x^*,\mu^*) &
            =0,\cdots, \frac{\partial L}{\partial
              \mu_m}(x^*,\mu^*)=0 \nonumber
	\end{align}	
	\end{block}
\end{frame}

	

	
	
	
	

	

\section{Inequality Constraints}
\begin{frame}
	\frametitle{Inequality Constraints}
\begin{itemize}
	\item Many optimization problems have their constraints defined by inequalities
	$g_1(x_1,\cdots,x_n)\leq b_1,\cdots,g_k(x_1,\cdots,x_n)\leq b_k$
	\item The method for finding the constrained maxima is more complex
	\item We start with the simplest case when there are two variables and
    one \emph{inequality} constraint
\end{itemize}
\end{frame}	

\begin{frame}
	\frametitle{Inequality Constraints}
	\begin{itemize}
		\item Consider the optimization problem with one inequality constraint:
		\begin{align}
		\max\  & f(x,y) \nonumber  \\
		\text{s.t.}\ \ & g(x,y)\leq b \nonumber 
		\end{align}
	\end{itemize}
\end{frame}	

\begin{frame}
	\frametitle{Inequality Constraints}
Case A: If the maximum of $f$ occurs at a point where $g(x,y)=b$, that is, when the constraint is \textbf{binding} (active/effective/tight)
	\begin{itemize}
		\item Highest level curve of $f$ meets the constraint set at point $\mathbf{p}$
		\item Equivalently, the level set of $f$ and the level set of $g$ are
        tangent to each other at $\mathbf{p}$
		\item $\triangledown f(\mathbf{p})$, $\triangledown g(\mathbf{p})$
        line up: $\triangledown f(\mathbf{p}) - \lambda \triangledown
        g(\mathbf{p}) = \mathbf{0}$
		\item More importantly, $\triangledown f(\mathbf{p})$ and $\triangledown g(\mathbf{p})$ \emph{point to the same direction}: {\color{red} $\lambda \geq 0 $}
		\item Form the Lagrangian function: $$L(x,y,\lambda) = f(x,y)-\lambda[g(x,y)-b]$$
	\end{itemize}
\end{frame}	

\begin{frame}
	\frametitle{Inequality Constraints}
	Case B: If the maximum of $f$ occurs at a point where $g(x,y)<b$, that is, when the constraint is \textbf{not binding} (inactive/ineffective/loose)
	\begin{itemize} 
		\item $\mathbf{q}$ must be a local max of $f$, that is, a local
        unconstrained max
		\item Then we have $\frac{\partial f}{\partial x}(\mathbf{q}) = 0$
        and $\frac{\partial f}{\partial y}(\mathbf{q}) = 0$
		\item The derivative of $g$ does not enter the calculations at $\mathbf{q}$ 
		\item In this case, we can still use the Lagrangian
        function: $$L(x,y,\lambda) = f(x,y)-\lambda[g(x,y)-b]$$ provided that
        we set {\color{red} $\lambda = 0$}.
	\end{itemize}
\end{frame}	

\begin{frame}
	\frametitle{Inequality Constraints}
In summary:
	\begin{itemize} 
		\item The constraint is binding: $g(x,y)-b = 0$, in this case $\lambda \geq 0$
		\item The constraint is not binding, in this case $\lambda = 0$
		\item Therefore $\lambda=0$ or $g(x,y)-b=0$ (and both may hold). This
        is called a \textbf{complementary slackness condition}:
        $$\lambda(g(x,y)-b) = 0$$
		\item Since we do not know if the constraint will be binding at the
        maximizer, we use the above equation to replace $\partial L/\partial
        \lambda = 0$
	\end{itemize}
\end{frame}	


\begin{frame}
	\frametitle{Inequality Constraints: Theorem}
    \begin{block}{\textbf{Theorem 18.3}}
        Suppose that $f$ and $g$ are $C^1$ and $(x^*,y^*)$ maximizes $f$ on
        the set $g(x,y)\leq b$. If $g(x^*,y^*) = b$, further suppose that
        $\frac{\partial g}{\partial x}(x^*,y^*)\neq 0$ or $\frac{\partial
          g}{\partial y}(x^*,y^*)\neq 0$. In any case, form the Lagrangian
        function
        $$L(x,y,\lambda) \equiv f(x,y)-\lambda(g(x,y)-b)$$
        Then there is a multiplier $\lambda^*$ such that
        \begin{itemize}
            \item (a) $\frac{\partial L}{\partial x}(x^*,y^*,\lambda^*) = 0$
            \item (b) $\frac{\partial L}{\partial y}(x^*,y^*,\lambda^*) = 0$
            \item (c) $\lambda^*[g(x^*,y^*)-b] = 0$ 
            \item (d) $\lambda^* \geq 0$ 
            \item (e) $g(x^*,y^*) \leq b$
        \end{itemize}
    \end{block}
\end{frame}	


\begin{frame}
\frametitle{Inequality Constraints: Example}
\textbf{Example 18.7}: Consider the problem: 
\begin{align}
\max\ \ f(x,y) & = xy \nonumber \\
\text{s.t.}\ \ \ g(x,y) &= x^2+y^2 \leq 1 \nonumber 
\end{align}
The only critical point of g occurs at the origin -- far away from the boundary of the constraint set $x^2+y^2 = 1$.  So the constraint qualification will
be satisfied at any candidate for a solution. Form the Lagrangian function $$L(x,y,\lambda) = xy-\lambda(x^2+y^2-1)$$ 
and write out the first order conditions described
in Theorem 18.3: 
\end{frame}	


\begin{frame}
	\frametitle{Inequality Constraints: Example}
	\textbf{Example 18.7} : 
	\begin{align}
	\frac{\partial L}{\partial x}  = y-2\lambda x & = 0 \nonumber \\
	\frac{\partial L}{\partial y}  = x-2\lambda y & = 0 \nonumber \\
	\lambda(x^2+y^2-1) & = 0 \nonumber \\
	x^2+y^2 & \leq 1 \nonumber \\
	\lambda & \geq 0 \nonumber
	\end{align}
The first two equations yield
$$\lambda = \frac{y}{2x} = \frac{x}{2y} \ \ \text{or}\ \ x^2 = y^2$$
\end{frame}	


\begin{frame}
	\frametitle{Inequality Constraints: Example}
	\textbf{Example 18.7} : 
\begin{itemize}
\item If $\lambda = 0$, then $x = y = 0$, which is a candidate for
a solution. 
\item If $\lambda \neq 0$, then $x^2+y^2-1 = 0$ $\Rightarrow $ $x^2 = y^2 =
\frac{1}{2}$, which gives $x = \pm \frac{1}{\sqrt{2}}, y = \pm
\frac{1}{\sqrt{2}}$ then we find the following four candidates: \small
\begin{align}
x & = +\frac{1}{\sqrt{2}}, y = +\frac{1}{\sqrt{2}}, \lambda = +\frac{1}{2} \nonumber \\
x & = -\frac{1}{\sqrt{2}}, y = -\frac{1}{\sqrt{2}}, \lambda = +\frac{1}{2} \nonumber \\
x & = +\frac{1}{\sqrt{2}}, y = -\frac{1}{\sqrt{2}}, \lambda = -\frac{1}{2} \nonumber \\
x & = -\frac{1}{\sqrt{2}}, y = +\frac{1}{\sqrt{2}}, \lambda = -\frac{1}{2} \nonumber 
\end{align} 
\end{itemize}
\end{frame}	

\begin{frame}
	\frametitle{Inequality Constraints: Example} \textbf{Example 18.7} : We
    disregard the last two candidates since they involve a negative
    multiplier. Plugging the three candidates into the object function, we
    find that
	$$x = \frac{1}{\sqrt{2}},\, y = \frac{1}{\sqrt{2}} \text{ and } x =
    -\frac{1}{\sqrt{2}},\, y = -\frac{1}{\sqrt{2}}$$
    are the solutions of our original problem. The two points with the
    negative multipliers are the solutions of the problem of minimizing $xy$
    on the constraint set $x^2+y^2\leq 1$.
\end{frame}




\begin{frame}
	\frametitle{Several Inequality Constraints}
\begin{itemize}
	\item Problem setup:
	\begin{align}
	\max \ \ &f(x) \nonumber \\
	\text{s.t.}\ \ &g_1(x) \leq b_1 \nonumber \\
    &g_2(x) \leq b_2 \nonumber \\
	 & \quad\vdots \nonumber \\
	& g_k(x) \leq b_k \nonumber
	\end{align}
\end{itemize}
\end{frame}


\begin{frame}
	\frametitle{Inequality Constraints: Theorem}
    \begin{block}{\textbf{Theorem 18.4}}
        Suppose $x^*$ is a local maximizer of $f$ on the constraint set.
        \begin{itemize}
            \item Assume the first $k_0$ constraints are binding at
            $x^*$ and the last $k-k_0$ are not binding.
            \item Suppose the following NDCQ condition is satisfied at
            $x^*$: the rank of
            $$\left({\begin{array}{ccc} 
                      \frac{\partial g_1}{\partial x_1}(x^*) & \cdots &\frac{\partial g_1}{\partial x_n}(x^*)\\
                      \vdots & \vdots & \vdots \\	
                      \frac{\partial g_{k_0}}{\partial x_1}(x^*) & \cdots &\frac{\partial g_{k_0}}{\partial x_n}(x^*)\\		
                  \end{array}}\right) $$ is $k_0$ -- as large as it can be.
        \end{itemize}		    
    \end{block}
    
\end{frame}	

\begin{frame}
    \frametitle{Inequality Constraints: Theorem}
    \begin{block}{\textbf{Theorem 18.4 (Cond.)}}
        Form the Lagrangian function: 
        \begin{equation*}
            L(x_1,\cdots,x_n,\lambda_1,\cdots,\lambda_k)  \equiv f(x) -
            \lambda_1[g_1(x)-b_1] - \ldots - \lambda_k[g_k(x)-b_k].
        \end{equation*}
        Then there exist multipliers $\lambda_1^*,\cdots, \lambda_k^*$ such that \small
        \begin{itemize}
            \item (a) $\frac{\partial L}{\partial x_1}(x^*,\lambda^*) = 0,\cdots, \frac{\partial L}{\partial x_n}(x^*,\lambda^*) = 0 $
            \item (b) $\lambda_1^*[g_1(x^*)-b_1] = 0, \cdots, \lambda_k^*[g_k(x^*)-b_k] = 0$ 
            \item (c) $\lambda_1^* \geq 0, \cdots, \lambda_k^* \geq 0$
            \item (d) $g_1(x^*)\leq b_1,\cdots, g_k(x^*)\leq b_k$
        \end{itemize}		
    \end{block}
\end{frame}	





\section{Mixed Constraints}

\begin{frame}
\frametitle{Mixed Constraints}
\begin{itemize}
	\item Some maximization problems involve both equality and inequality constraints
	\item Problem setup:
		\begin{align}
		\max\  & f(\mathbf{x}) \nonumber  \\
		\text{s.t.}\ \ & g_1(\mathbf{x})\leq b_1,\cdots, g_k(\mathbf{x})\leq b_k   \nonumber \\
		& h_1(\mathbf{x}) = c_1,\cdots, h_m(\mathbf{x})= c_m \nonumber 
		\end{align}
\end{itemize}
\end{frame}

\begin{frame}
	\frametitle{Mixed Constraints: Theorem}
    \begin{block}{\textbf{Theorem 18.5}}
        Suppose that $f$, $g_1,\cdots,g_k$, $h_1,\cdots,h_m$ are $C^1$.
        Suppose that $x^*\in \mathbb{R}^n$ is a local maximizer of
        $f$ on the constraint set.
        \begin{itemize}
            \item Assume the first $k_0$ constraints are binding at
            $x^*$ and the last $k-k_0$ are not binding. Suppose the
            following NDCQ condition is satisfied at $x^*$: the rank of
            $$\left({\begin{array}{ccc} 
                      \frac{\partial g_1}{\partial x_1}(x^*) & \cdots &\frac{\partial g_1}{\partial x_n}(x^*)\\
                      \vdots & \vdots & \vdots \\	
                      \frac{\partial g_{k_0}}{\partial x_1}(x^*) & \cdots &\frac{\partial g_{k_0}}{\partial x_n}(x^*)\\
                      \frac{\partial h_1}{\partial x_1}(x^*) & \cdots &\frac{\partial h_1}{\partial x_n}(x^*)\\	
                      \vdots & \ddots & \vdots \\
                      \frac{\partial h_m}{\partial x_1}(x^*) & \cdots &\frac{\partial h_m}{\partial x_n}(x^*)\\	
                  \end{array}}\right)  $$
            is $k_0+m$ -- as large as it can be. 
        \end{itemize}    
    \end{block}
	
\end{frame}	



\begin{frame}
	\frametitle{Mixed Constraints: Theorem}
    \begin{block}{\textbf{Theorem 18.5 (Cond.)}}
        Form the Lagrangian Function \small
        \begin{align} 
            L(x_1,\cdots,x_n,\lambda_1,\cdots,\lambda_k,\mu_1,\cdots,\mu_m)
            \equiv f(x)&-\lambda_1[g_1(x)-b_1]-\cdots - \lambda_k[g_k(x)-b_k] \nonumber \\
            & -\mu_1[h_1(x)-c_1]-\cdots - \mu_m[h_m(x)-c_m]
            \nonumber
        \end{align} 
        Then there exist multipliers $\lambda_1^*,\cdots,
        \lambda_k^*,\mu_1^*,\cdots,\mu_m^*$ such that:
        \begin{itemize}
            \item (a) $\frac{\partial L}{\partial x_1}(x^*,\lambda^*,\mu^*) = 0,\cdots, \frac{\partial L}{\partial x_n}(x^*,\lambda^*,\mu^*) = 0$
            \item (b) $\lambda_1^*[g_1(x^*)-b_1] = 0,\cdots, \lambda_k^*[g_k(x^*)-b_k] = 0$ 
            \item (c) $h_1(x^*) = c_1,\cdots,h_m(x^*) = c_m$ 
            \item (d) $\lambda_1^* \geq 0,\cdots, \lambda_k^* \geq 0$ 
            \item (e) $g_1(x^*) \leq b_1,\cdots, g_k(x^*) \leq b_k$
        \end{itemize}
    \end{block}
\end{frame}	

% \begin{frame}
% 	\frametitle{Exercise}
% 	\textbf{Exercise} Suppose we can buy a chemical for \$10 per ounce. There
%     are only 17.25 oz available. We can transform this chemical into two
%     products: A and B (assume one oz of this chemical can be transformed
%     equivalently into one oz of the products). Transforming to A costs \$3
%     per oz, while transforming to B costs \$5 per oz. If $x_1$ oz of A are
%     produced, the price we command for A is \$$30-x_1$; if $x_2$ oz of B are
%     produced, the price we get for B is \$$50-x_2$. How much chemical should
%     we buy? and what should we transform it to?
% \end{frame}

%
%\begin{frame}
%\frametitle{Homework}
%\begin{itemize}
%	\item Exercise 18.7
%	\item Exercise 18.12
%	\item Exercise 18.17
%\end{itemize}
%\end{frame}


\end{document}
